\section*{अध्याय \ref{ch:b2:batteries-electrolysis} --- बैटरियाँ, ईंधन सेल और विद्युत-अपघटन}

\begin{solution}{exo:b2:batteries-electrolysis:1}
अकेला ज़िंक: लगभग \qty{-0.82}{V} और 0.15 की धारा (मॉडल इकाइयाँ)। ताँबे को
छूता हुआ: लगभग \qty{-0.74}{V} और 2.4। धारा, और इसलिए विलयन की दर,
ताँबे के संपर्क में लगभग सोलह गुना अधिक है।
\end{solution}

\begin{solution}{exo:b2:batteries-electrolysis:2}
$m = 500 \times 3600 \times 65.38/(2 \times 96\,485) = \qty{610}{g}$।
\end{solution}

\begin{solution}{exo:b2:batteries-electrolysis:3}
$E^\circ = \qty{2.04}{V}$ (अध्याय की तरह); श्रेणी में छह सेल
\qty{12}{V} देते हैं।
\end{solution}

\begin{solution}{exo:b2:batteries-electrolysis:4}
$1/6.94 = \qty{0.144}{mol}$; $0.144 \times 96\,485 = \qty{1.39e4}{C}$, अर्थात्
\qty{3.86}{A.h}।
\end{solution}

\begin{solution}{exo:b2:batteries-electrolysis:5}
100~\% पर: $200 \times 86\,400 \times 63.546/(2 \times 96\,485) = \qty{5.69}{kg}$;
$\eta_F = 5.50/5.69 = 0.967$।
\end{solution}

\begin{solution}{exo:b2:batteries-electrolysis:6}
$w = 2 \times 96\,485 \times 3.3/(0.90 \times 0.06538) = \qty{1.08e7}{J/kg}$, अर्थात्
\qty{3.0}{kW.h/kg}।
\end{solution}

\begin{solution}{exo:b2:batteries-electrolysis:7}
$U = E - rI$: $r = (1.55 - 1.40)/(0.60 - 0.10) = \qty{0.30}{\ohm}$; $E = 1.55 + 0.30
\times 0.10 = \qty{1.58}{V}$।
\end{solution}

\begin{solution}{exo:b2:batteries-electrolysis:8}
ऐनोड \ce{2Cl- -> Cl2 + 2e-}; कैथोड \ce{2H2O + 2e- -> H2 + 2OH-}। $\Delta_r G^\circ =
2(-157.24) - 2(-131.23) - 2(-237.13) = \qty{422.2}{kJ/mol}$, $U_{\min} =
\qty{2.19}{V}$। क्लोरीन हाइड्रॉक्साइड से अभिक्रिया करती (\ce{Cl2 + 2OH- -> Cl- +
ClO- + H2O}), जिससे दोनों उत्पाद बिगड़ जाते, और क्लोरीन तथा हाइड्रोजन का मिश्रण
विस्फोट कर सकता है।
\end{solution}

\begin{solution}{exo:b2:batteries-electrolysis:9}
\ce{NaCl -> Na + 1/2Cl2}, $n = 1$: $U_{\min} = 310\,674/96\,485 = \qty{3.22}{V}$।
प्रति किलोग्राम सोडियम: $96\,485 \times 3.22/0.02299 = \qty{1.35e7}{J}$, अर्थात्
\qty{3.75}{kW.h}।
\end{solution}

\begin{solution}{exo:b2:batteries-electrolysis:10}
ज़िंक कैथोड पर \ce{H+} का अपचयन मंद है और उसकी भित्ति ज़िंक
तरंग से नीचे है: विभव घटाने पर पहले ज़िंक जमता है, थोड़ी
हाइड्रोजन के साथ। प्लैटिनम पर \ce{H+} का अपचयन तीव्र है और लगभग
\qty{0}{V} पर आरंभ होता है: आरंभ में केवल हाइड्रोजन निकलती है। प्लैटिनम के
ज़िंक से ढक जाने पर कैथोड ज़िंक कैथोड जैसा व्यवहार करता है और ज़िंक जमता है।
\end{solution}

\begin{solution}{exo:b2:batteries-electrolysis:11}
\qty{1250}{K} पर: $\Delta_f G^\circ(\ce{Al2O3}) = -1328.3 + 0.75 \times 66.0 =
\qty{-1278.8}{kJ/mol}$, $\Delta_f G^\circ(\ce{CO2}) = \qty{-396.1}{kJ/mol}$। $\Delta_r G^\circ =
3(-396.1) - 2(-1278.8) = \qty{1369}{kJ}$, $n = 12$: $U_{\min} = \qty{1.18}{V}$। अक्रिय ऐनोड:
$\Delta_r G^\circ = \qty{2557.5}{kJ}$, $U_{\min} = \qty{2.21}{V}$; जलता कार्बन
कार्य का लगभग \qty{1.03}{V} प्रदान करता है।
\end{solution}

\begin{solution}{exo:b2:batteries-electrolysis:12}
प्रति मोल जल आवेश $2F$ दोनों दिशाओं में समान है: दक्षता $0.70/1.80 =
0.39$। शेष ऊष्मा है: दोनों ऑक्सीजन इलेक्ट्रोडों के \omterm{def:b2:current-potential-curves:overpotential}{अधिविभव} (दोनों
दिशाओं में मंद), हाइड्रोजन इलेक्ट्रोडों के \omterm{def:b2:current-potential-curves:overpotential}{अधिविभव}, और ओमीय हानियाँ।
\end{solution}

\begin{solution}{pb:b2:batteries-electrolysis:1}
\textbf{1.} ऐलुमिनियम $+3 \to 0$; ऑक्सीजन $-2$ ही रहती है; कार्बन $0 \to +4$।
\textbf{2.} कैथोड \ce{Al^3+ + 3e- -> Al}; ऐनोड \ce{C + 2O^2- -> CO2 + 4e-}।
\textbf{3.} \ce{2Al2O3 + 3C -> 4Al + 3CO2}, $n = 12$।
\textbf{4.} \ce{Al2O3}: $-1328.3 + 0.75 \times (1328.3 - 1262.3) = \qty{-1278.8}{kJ/mol}$;
\ce{CO2}: \qty{-396.1}{kJ/mol}।
\textbf{5.} $3(-396.1) - 2(-1278.8) = \qty{+1369}{kJ}$।
\textbf{6.} $1\,369\,000/(12 \times 96\,485) = \qty{1.18}{V}$।
\textbf{7.} $2 \times 1278.8/(12 \times 96.485) = \qty{2.21}{V}$: कार्बन का ऑक्सीकरण,
जो स्वयं \omterm{def:b2:reaction-free-energy:exergonic}{ऊर्जामोची} है, ऑक्साइड के विघटन के कार्य का लगभग आधा चुकाता है।
\textbf{8.} $10^6/26.982 = \qty{3.706e4}{mol}$।
\textbf{9.} $3 \times 96\,485 \times 3.706 \times 10^4 = \qty{1.073e10}{C}$।
\textbf{10.} $1.073 \times 10^{10}/0.94 = \qty{1.141e10}{C}$।
\textbf{11.} $1.141 \times 10^{10}/3.0 \times 10^5 = \qty{3.80e4}{s}$, \qty{10.6}{h}।
\textbf{12.} $24/10.6 = 2.27$: प्रति दिन \qty{2.27}{t}।
\textbf{13.} कुंड में घुले ऐलुमिनियम का एक भाग ऐनोड गैस तक पहुँचकर
कार्बन डाइऑक्साइड द्वारा फिर ऑक्सीकृत हो जाता है, इसलिए उसका आवेश दो बार खर्च होता है; कुछ धारा
बिना \omterm{def:b2:batteries-electrolysis:electrolysis}{विद्युत-अपघटन} किए कुंड से रिस भी जाती है।
\textbf{14.} $1.141 \times 10^{10} \times 4.1 = \qty{4.68e10}{J}$, \qty{13.0}{MW.h}।
\textbf{15.} \qty{13.0}{kW.h/kg}।
\textbf{16.} $3 \times 96\,485 \times 1.18/0.026982 = \qty{1.27e7}{J/kg}$, \qty{3.5}{kW.h/kg}।
\textbf{17.} ऊष्मा में: इलेक्ट्रोडों के \omterm{def:b2:current-potential-curves:overpotential}{अधिविभव} और, मुख्यतः, कुंड में
ओमीय पात। वह ऊष्मा कुंड को लगभग
\qty{1250}{K} पर पिघला रखती है, जिसके लिए अन्यथा ईंधन चाहिए होता।
\textbf{18.} $4.1 \times 3.0 \times 10^5 = \qty{1.23}{MW}$।
\textbf{19.} $28\,690/0.026982 = \qty{1.06e6}{J/kg}$, \qty{0.30}{kW.h/kg}: \omterm{def:b2:batteries-electrolysis:electrolysis}{विद्युत-अपघटन} की
ऊर्जा का लगभग 2~\%।
\textbf{20.} $1.18/4.1 = 0.29$।
\textbf{21.} $3.706 \times 10^4 \times 3/4 \times 12.011 = \qty{334}{kg}$।
\textbf{22.} $3.706 \times 10^4 \times 3/4 \times 44.009 = \qty{1.22e3}{kg}$।
\textbf{23.} ऐनोडों के गर्म ऊपरी भाग वायु में जलते हैं; कार्बन का एक भाग
कार्बन मोनॉक्साइड के रूप में निकलता है (उस ताप पर \ce{C + CO2 -> 2CO}), जो प्रति ऑक्सीजन परमाणु
दुगुना कार्बन लेती है।
\textbf{24.} प्रति किलोग्राम ऐलुमिनियम \ce{CO2} के \qty{1.22}{kg}।
\textbf{25.} न्यूनतम वोल्टता लगभग \qty{1.03}{V} बढ़ जाती, परंतु ऐनोड
कार्बन डाइऑक्साइड के बजाय ऑक्सीजन छोड़ता और उपभुक्त नहीं होता।
\textbf{26.} सेल प्रति किलोग्राम ऐलुमिनियम $\boldsymbol{\approx \qty{13}{kW.h}}$ विद्युत
ऊर्जा प्रयुक्त करता है।
\end{solution}
